(1) 粒子质量 \(m\) 在边长 \(L\) 的三维无限立方势阱 \(0<x,y,z<L\) 中运动，求归一能量本征函数。
(2) 求能谱，并列出基态及第一激发能级的简并度。
(3) 计算基态的能量和位置期望 \(\langle x\rangle\) 与 \(\langle x^2\rangle\)。
(4) 在高能大量子数下用正八分球计数，求能量不超过 \(E\) 的态数渐近标度。
